\subsection*{프로젝트 V-모델 및 추적 구조}

본 문서는 주행 상황 실시간 경고 시스템 프로젝트의 개발 문서와 검증 문서가 실제 V-모델 구조 안에서 어떻게 대응되는지를 reviewer-facing 형식으로 정리한 별첨이다.

본 프로젝트는 CANoe SIL 기반 구현형 프로젝트로, ASPICE 전체 산출물을 모두 독립 문서로 분리하지는 않았다. 대신 \texttt{01 \textrightarrow{} 03 \textrightarrow{} 0301 \textrightarrow{} 0302 \textrightarrow{} 0303 \textrightarrow{} 0304 \textrightarrow{} 04 \textrightarrow{} 05 \textrightarrow{} 06 \textrightarrow{} 07} 문서 체인을 실행 가능한 기준선으로 유지하고, \texttt{00d\_HARA\_Worksheet}와 \texttt{00g\_Master\_Test\_Matrix}를 통해 안전성 및 추적성을 고정하였다.

\paragraph{프로젝트 V-모델 해석}
\begin{itemize}
\item 좌측 개발 축은 요구사항, 시스템 설계, 통신/변수 계약, CAPL 구현 순으로 구성하였다.
\item 우측 검증 축은 단위시험, 통합시험, 시스템시험 순으로 구성하였다.
\item \texttt{00d\_HARA\_Worksheet}는 안전 사건과 검증 우선순위의 기준점으로 사용하였다.
\item \texttt{00g\_Master\_Test\_Matrix}는 요구사항, 기능, 시험 자산을 연결하는 프로젝트 Test Matrix로 사용하였다.
\end{itemize}

\paragraph{V-모델 매핑 표}
\renewcommand{\arraystretch}{1.15}
\begin{tabularx}{\textwidth}{>{\raggedright\arraybackslash}p{0.17\textwidth} >{\raggedright\arraybackslash}p{0.20\textwidth} >{\raggedright\arraybackslash}X >{\raggedright\arraybackslash}p{0.18\textwidth}}
\toprule
관점 & ASPICE / ISO 26262 & 프로젝트 문서 & 검증 대응 \\
\midrule
안전 / 추적성 기준 & SYS.2, SUP.10 / Part 3, Part 8 & 00d\_HARA\_Worksheet, 00g\_Master\_Test\_Matrix & 05, 06, 07 전 계층 \\
시스템 요구사항 & SYS.2 / Part 4, Cl.6 & 01\_Requirements & 07\_System\_Test \\
시스템 설계 & SYS.3 / Part 4, Cl.7 & 03, 0301, 0302 & 06, 07 \\
통신 / 변수 계약 & SWE.2, SWE.3 / Part 6, Cl.7--8 & 0303, 0304 & 05, 06 \\
구현 기준선 & SWE.3 / Part 6, Cl.9 & 04\_SW\_Implementation & 05, 06 \\
단위 검증 & SWE.4 / Part 6, Cl.9 & 05\_Unit\_Test & 0303, 0304, 04 \\
통합 검증 & SWE.5 / Part 6, Cl.10 & 06\_Integration\_Test & 03, 0301, 0302, 0303, 0304 \\
시스템 검증 & SYS.5 / Part 4, Cl.10 & 07\_System\_Test & 01\_Requirements, 00d\_HARA\_Worksheet \\
\bottomrule
\end{tabularx}

\paragraph{표의 해석}
\begin{itemize}
\item \texttt{00d\_HARA\_Worksheet}는 위험 사건과 안전 우선순위의 기준점이다.
\item \texttt{00g\_Master\_Test\_Matrix}는 요구사항, 기능, 시험 자산을 연결하는 프로젝트 Test Matrix이다.
\item \texttt{03}, \texttt{0301}, \texttt{0302}는 ECU 역할, 기능 분해, 네트워크 흐름을 담당하는 시스템 설계 축이다.
\item \texttt{0303}, \texttt{0304}는 실행 가능한 통신/변수 계약을 정의하는 인터페이스 축이다.
\item \texttt{04\_SW\_Implementation}은 CAPL, panel, sysvar를 포함한 구현 기준선이다.
\end{itemize}

\paragraph{프로젝트 추적성 체인}
본 프로젝트의 추적성 체인은 아래 순서로 고정하였다.

\begin{center}
\texttt{Req \textrightarrow{} Func \textrightarrow{} Flow \textrightarrow{} Comm \textrightarrow{} Var \textrightarrow{} Code \textrightarrow{} Unit Test \textrightarrow{} Integration Test \textrightarrow{} System Test}
\end{center}

\paragraph{작성 기준}
\begin{itemize}
\item 본 매핑은 현재 active baseline 문서를 기준으로 정리하였다.
\item 별도 SWE.1 독립 문서는 두지 않고, 기능 정의와 통신/변수 계약 문서가 소프트웨어 요구/설계 역할을 분담하도록 구성하였다.
\item 따라서 본 프로젝트의 V-모델은 이론적 완전 분리형이 아니라, CANoe SIL 구현 프로젝트에 맞춘 실행형 추적 구조로 해석한다.
\end{itemize}
